\documentclass[11pt]{article}
\usepackage[margin=2.1cm]{geometry}
\usepackage{booktabs}
\usepackage{amsmath}
\usepackage{amssymb}
\usepackage{enumitem}
\usepackage{array}
\usepackage{listings}
\usepackage{xcolor}
\usepackage[T1]{fontenc}
\usepackage{lmodern}

\setlength{\parindent}{0pt}
\setlength{\parskip}{3.6pt}
\newcommand{\code}[1]{\texttt{#1}}
\newcommand{\hr}{\vspace{2pt}\hrule\vspace{6pt}}
\renewcommand{\arraystretch}{1.12}

\definecolor{kw}{RGB}{0,90,160}
\definecolor{cmt}{RGB}{110,110,110}
\lstdefinestyle{vhdl}{
  language=VHDL, basicstyle=\ttfamily\footnotesize,
  keywordstyle=\color{kw}\bfseries, commentstyle=\color{cmt}\itshape,
  showstringspaces=false, columns=fullflexible, breaklines=true,
  frame=leftline, framesep=6pt, xleftmargin=8pt, aboveskip=6pt, belowskip=6pt
}
\lstset{style=vhdl}

\begin{document}

\begin{center}
{\Large\bfseries EE\,362 --- Pre-Lab 8: Solution}\\[2pt]
{\large Next Instruction Address Generation II}\\[4pt]
Fall 2026 \quad$\cdot$\quad 5 points \quad$\cdot$\quad TA marking guide
\end{center}
\hr

{\small Marking: Q1 = 2 points, Q2 = 2 points, Q3 = 1 point.}

\vspace{4pt}\hr

\section*{Question 1 --- The six equations (2 points)}

\subsection*{(a) Taken, in terms of $Z$ and $N$}

\begin{center}
\begin{tabular}{lll}
\toprule
code & meaning & Taken $=$ \\
\midrule
\code{NONE} & always & $1$ \\
\code{EQ} & $= 0$ & $Z$ \\
\code{NEQ} & $\neq 0$ & $\overline{Z}$ \\
\code{GTZ} & $> 0$ & $\overline{N} \cdot \overline{Z}$ \\
\code{GEZ} & $\geq 0$ & $\overline{N}$ \\
\code{LTZ} & $< 0$ & $N$ \\
\code{LEZ} & $\leq 0$ & $N + Z$ \\
\bottomrule
\end{tabular}
\end{center}

\textbf{The two that carry the marks are \code{GTZ} and \code{LEZ}}, because they
are the only ones needing \emph{both} flags. ``Greater than zero'' means not
negative \textbf{and} not zero; ``less than or equal'' means negative
\textbf{or} zero.

\subsection*{(b) The three complement pairs}

\[
\code{GTZ} \leftrightarrow \code{LEZ}
\qquad
\code{GEZ} \leftrightarrow \code{LTZ}
\qquad
\code{EQ} \leftrightarrow \code{NEQ}
\]

\begin{quote}
\textbf{Model answer:} Each pair is one inverter apart, so the seven conditions do
not cost seven circuits. Three small gates and their complements cover all of
them, and the condition code just selects which one reaches the output ---
\textbf{the whole condition unit is a handful of gates and a multiplexer.}
\end{quote}

Check the complement algebraically once at the board:
$\overline{\code{GTZ}} = \overline{\overline{N}\cdot\overline{Z}} = N + Z = \code{LEZ}$
by De Morgan's law. That is the kind of line their report should contain.

\subsection*{(c) A classmate writes \code{GTZ} as $\overline{N}$}

\begin{center}
\begin{tabular}{ccccl}
\toprule
$N$ & $Z$ & correct \code{GTZ} & their version & \\
\midrule
0 & 0 & 1 & 1 & agree (result positive) \\
\textbf{0} & \textbf{1} & \textbf{0} & \textbf{1} & \textbf{disagree --- result is exactly zero} \\
1 & 0 & 0 & 0 & agree (result negative) \\
1 & 1 & 0 & 0 & agree (cannot occur) \\
\bottomrule
\end{tabular}
\end{center}

\begin{quote}
\textbf{Model answer:} They disagree only at $N = 0$, $Z = 1$ --- a result of
exactly zero. Their version is $\overline{N}$, which is \code{GEZ}, so a
\code{BGTZ} would branch when the value is zero. A loop that should stop at zero
would take the branch one extra time.
\end{quote}

\textbf{Mark strictly:} full credit needs the specific flag pair $N{=}0, Z{=}1$,
not just ``when the result is zero''. The same slip in \code{LEZ} (writing $N$ and
forgetting $+Z$) is the one that actually bites in Lab 12, because the reference
program's loop ends with \code{BLEZ}.

\vspace{4pt}\hr

\section*{Question 2 --- Absolute jumps (2 points)}

\subsection*{(a) \code{J 0x08} at \code{0x00000014}, encoding \code{0x08000002}}

\begin{center}
\begin{tabular}{lll}
\toprule
piece & value & working \\
\midrule
\code{NSA} & \code{0x00000018} & $\code{0x14} + 4$ \\
top 4 bits & \code{0000} & inherited from \code{NSA(31..28)} \\
\code{Index} & \code{0x0000002} & the low 26 bits of the encoding \\
\textbf{\code{Abslt}} & \textbf{\code{0x00000008}} & \code{NSA(31..28)} \& \code{Index} \& \code{"00"} $= 2 \times 4$ \\
\bottomrule
\end{tabular}
\end{center}

\textbf{No arithmetic happens here.} \code{Abslt} is pure concatenation --- wiring,
not an adder. A student who wrote an addition has built hardware the machine does
not need.

\subsection*{(b) Where the 32 bits come from, and the range limit}

\begin{center}
\begin{tabular}{lll}
\toprule
bits & from & why \\
\midrule
4 & \code{NSA(31..28)} & \textbf{inherited} from where we already are \\
26 & the instruction's \code{Index} field & all the room the encoding has \\
2 & constant \code{"00"} & instructions are word-aligned, so these are always zero \\
\bottomrule
\end{tabular}
\end{center}

\begin{quote}
\textbf{The range limit:} the top four bits cannot be changed by the instruction,
so a \code{J} can only reach addresses inside the \textbf{same 256\,MB region}
($2^{28}$ bytes) as the instruction after it. Jumping across a region boundary is
impossible with \code{J}.
\end{quote}

Full credit needs $2^{28}$ (or 256\,MB) \emph{and} the reason --- the inherited
top bits.

\subsection*{(c) When you need to go further}

\begin{quote}
\textbf{Model answer:} A jump through a register (\code{JR}). The address comes
from a 32-bit register instead of from the instruction, so all 32 bits are
supplied at run time and nothing is inherited. That is why the \code{REG} path
exists in the multiplexer.
\end{quote}

\vspace{4pt}\hr

\section*{Question 3 --- Test plan (1 point)}

\subsection*{(a) Two pairs differing only in \code{Taken}}

\begin{center}
\small
\begin{tabular}{l>{\raggedright\arraybackslash}p{5.0cm}>{\raggedright\arraybackslash}p{6.4cm}}
\toprule
& vectors & expected \\
\midrule
\textbf{different} addresses & \code{Sel} = \code{REL}, \code{Taken} = 0 then 1 & \code{NSA}, then \code{Rel} --- \textbf{only the relative path is conditional} \\
\textbf{same} address & \code{Sel} = \code{ABS} (or \code{REG}, or \code{NSA}), \code{Taken} = 0 then 1 & \code{Abslt} both times --- \textbf{jumps are unconditional by construction} \\
\bottomrule
\end{tabular}
\end{center}

\begin{quote}
\textbf{The point:} \code{Sel} and \code{Cond} are two separate questions ---
\emph{which} address is the candidate, and \emph{is the relative branch taken}. A
student who merged them into one control signal cannot write the second pair, and
that is exactly the confusion the question is designed to surface.
\end{quote}

\textbf{And note why there is no fifth select code for ``branch not taken'':} a
branch that is not taken simply continues to the next instruction, which is
already \code{NSA}. The fall-back is built in.

\subsection*{(b) The scripted sequence}

A model answer walks the PC through every path in order:

\begin{center}
\small
\begin{tabular}{c>{\raggedright\arraybackslash}p{5.6cm}ll}
\toprule
\# & what happens & expected PC & what it proves \\
\midrule
1 & reset & \code{0x00000000} & the reset vector \\
2 & \code{Sel}=\code{NSA}, one edge & \code{0x00000004} & sequential execution \\
3 & \code{Sel}=\code{NSA}, one edge & \code{0x00000008} & it keeps advancing by 4 \\
4 & \code{Sel}=\code{REL}, \code{Taken}=0 & \code{0x0000000C} & not-taken falls back to \code{NSA} \\
5 & \code{Sel}=\code{REL}, \code{Taken}=1, offset $+2$ & target $=$ \code{NSA}$+8$ & the taken branch \\
6 & \code{Sel}=\code{REL}, \code{Taken}=1, offset $-4$ & target below \code{NSA} & \textbf{backward branches} \\
7 & \code{Sel}=\code{ABS} & \code{Abslt} & the jump path, unconditional \\
8 & \code{Sel}=\code{REG} & the register value & the \code{JR} path \\
\bottomrule
\end{tabular}
\end{center}

\begin{quote}
\textbf{Why a scripted sequence rather than isolated vectors:} the PC is stateful,
so each step's expected value depends on the one before. Writing the expected
address next to every step turns the testbench into a small program trace --- and
that is exactly what you will read in Lab 12 when the CPU runs.
\end{quote}

\vspace{6pt}
\hr
{\small \textbf{At the bench:} the two answers that predict trouble are a
\code{GTZ} written as $\overline{N}$, and a test plan with no backward branch.
Both produce a NIAG that passes its own tests and a CPU whose loop never ends.}

\end{document}
